\documentclass[12pt,a4paper]{article}
\usepackage[utf8]{inputenc}
\usepackage{graphicx}
\usepackage{amsmath}
\usepackage{amssymb}
\usepackage[colorlinks=true, linkcolor=black, urlcolor=black, citecolor=black]{hyperref}
\usepackage{booktabs}
\usepackage{float}
\usepackage{caption}
\usepackage{subcaption}
\usepackage{xcolor}
\usepackage{listings}
\usepackage{algorithm}
\usepackage{algpseudocode}
\usepackage{geometry}
\usepackage{lipsum}  
\usepackage{cmbright}
\usepackage{titlesec}
\usepackage{tabularx} 
\titlespacing*{\section}{0pt}{1ex plus 0.5ex minus .2ex}{1ex plus 0.3ex}
\titlespacing*{\subsection}{0pt}{0.8ex plus 0.4ex minus .2ex}{0.8ex plus 0.3ex}
\titlespacing*{\subsubsection}{0pt}{0.6ex plus 0.3ex minus .2ex}{0.6ex plus 0.2ex}


% Optional: Set how deep TOC and numbering should go
\setcounter{tocdepth}{3}
\setcounter{secnumdepth}{3}

\title{\textbf{Object Detection Models \\ YOLOv8, Faster R-CNN, DETR}}


\makeatletter
\def\@maketitle{%
  \newpage
  \null
  \vskip 1em%
  \begin{center}%
  \let \footnote \thanks
    {\LARGE \@title \par}%
    \vskip 1em%
  \end{center}%
  \par
  \vskip 1em}
\makeatother

\begin{document}

\maketitle

\vspace{12cm}

\begin{table}[H]
    \centering
    \renewcommand{\arraystretch}{1.5} % Adjust row height (1.0 is default)
    \setlength{\tabcolsep}{20pt}      % Adjust column width (default is 6pt)
    \begin{tabularx}{0.8\textwidth}{|X|c|}
        \hline
        \textbf{Name} & \textbf{ID} \\
        \hline
        Hour Mohamed   & 8119 \\
        Nadine Mohamed Tamish & 8092 \\
        Seif Maged     & 8387 \\
        \hline
    \end{tabularx}
\end{table}

\newpage
\tableofcontents
\newpage

\section{Introduction}

Object detection is a core problem in computer vision that focuses on identifying and locating objects within an image. It plays a vital role in many real-world applications such as autonomous driving, surveillance, medical imaging, and robotics.
Over time, various model architectures have been developed to improve the speed and accuracy of detection. These models are typically grouped into three main categories:
\begin{itemize}
    \item \textbf{One-stage detectors}, such as YOLOv8, which perform detection in a single step and are known for their real-time performance.
    \item \textbf{Two-stage detectors}, such as Faster R-CNN, which first propose object regions and then refine the results—typically achieving higher accuracy.
    \item \textbf{Transformer-based detectors}, like DETR, which introduce attention mechanisms to model global relationships in the image and simplify the detection pipeline.
\end{itemize}

This report presents a comparative analysis of these three representative models, focusing on their internal architectures, detection capabilities, and performance on benchmark datasets.
\section{Datasets}

\subsection{COCO Dataset}
 \textbf{(COCO)} dataset is a widely adopted benchmark for assessing object detection, instance segmentation, and image captioning algorithms. In our study, we utilize the 2017 edition, which comprises:
\begin{itemize}
    \item \textbf{Training images:} 118,287
    \item \textbf{Validation images:} 5,000
    \item \textbf{Evaluated categories:} 80 object classes
\end{itemize}

\textit{Note:} While COCO annotations list 91 classes in total, only 80 of them are officially evaluated. The remaining 11 classes are included in the category index but lack labeled instances within the main dataset splits.

\subsection{Pascal VOC 2007 Dataset}
The \textbf{Pascal Visual Object Classes (VOC) 2007} dataset is another key benchmark commonly used to evaluate the performance of object detection models. In this project, we rely on:
\begin{itemize}
    \item \textbf{Trainval set:} 5,011 annotated images
    \item \textbf{Test set:} 4,952 annotated images
    \item \textbf{Categories:} 20 distinct object classes
\end{itemize}

VOC 2007 offers a practical and compact dataset with varied object appearances, occlusions, and real-world imaging conditions. Its size and structure make it ideal for model prototyping and performance comparison across different architectures.

\section{Architectures}
\subsection{YOLOV8}
\begin{itemize}
\item \textbf {Backbone: } 
    \begin{itemize}
        \item \textbf {CBS Blocks (Conv-BN-SiLU):} Basic building block: Convolution → BatchNorm → SiLU (Swish) activation.
        \item \textbf {C2f Modules:} 
            \begin{itemize}
                \item Similar to CSP (Cross Stage Partial) modules, but more efficient. \item Split input, process one part with bottlenecks, then fuse back. 
            \end{itemize}
                
        \item \textbf {SPPF (Spatial Pyramid Pooling – Fast):}
            \begin{itemize}
                \item Helps to capture multi-scale features with minimal computation.
                \item Replaces traditional SPP with a faster variant.
            \end{itemize}
        \item \textbf {Downsampling via Strided Convolutions or MaxPool:}
            \begin{itemize}
                \item Reduces spatial resolution while increasing channel depth.
            \end{itemize}
    \end{itemize}
\item \textbf {Neck (feature extractor):} 
             Merges feature maps from different stages of the backbone to capture information at various scales.

\item \textbf{Head:} Responsible for making predictions. 
\end{itemize}
\subsection{Faster R-CNN}
\begin{itemize}
    \item \textbf{Backbone (Feature Extractor):}
    \begin{itemize}
        \item Typically a pre-trained CNN like ResNet-50 or ResNet-101.
        \item Extracts convolutional feature maps from the input image.
        \item Often includes Feature Pyramid Networks (FPN) to capture multi-scale features.
    \end{itemize}

    \item \textbf{Region Proposal Network (RPN):}
    \begin{itemize}
        \item A shallow network that slides over the feature map produced by the backbone.
        \item For each location, it predicts:
        \begin{itemize}
            \item \textbf{Objectness score:} whether the region contains an object.
            \item \textbf{Bounding box offsets:} refinements to default anchor boxes.
        \end{itemize}
        \item Outputs a set of region proposals (ROIs) with high objectness scores.
    \end{itemize}

    \item \textbf{Anchor Boxes:}
    \begin{itemize}
        \item Predefined boxes of different scales and aspect ratios used by the RPN.
        \item Anchor boxes are refined using regression offsets predicted by the RPN.
    \end{itemize}

    \item \textbf{ROI Align (Region of Interest Alignment):}
    \begin{itemize}
        \item Takes variable-sized region proposals and extracts fixed-size (e.g., 7×7) feature maps from the backbone feature map.
        \item ROI Align improves upon ROI Pooling by preserving spatial alignment (no rounding).
    \end{itemize}

    \item \textbf{Head (Two Fully Connected Layers):}
    \begin{itemize}
        \item Processes the aligned ROIs with two fully connected (FC) layers.
        \item Outputs two branches:
        \begin{itemize}
            \item \textbf{Classification Branch:}
            \begin{itemize}
                \item Uses a \textbf{Softmax} function.
                \item Outputs class probabilities (including background).
            \end{itemize}
            \item \textbf{Bounding Box Regression Branch:}
            \begin{itemize}
                \item Outputs four coordinates for each class: \texttt{[tx, ty, tw, th]}.
                \item These are used to refine the final bounding box locations.
            \end{itemize}
        \end{itemize}
    \end{itemize}

    \item \textbf{Post-processing:}
    \begin{itemize}
        \item Applies \textbf{Non-Maximum Suppression (NMS)} to remove redundant overlapping boxes.
        \item Final output is a set of object detections with class labels, confidence scores, and bounding boxes.
    \end{itemize}

    \item \textbf{Training Strategy:}
    \begin{itemize}
        \item Multi-task loss = classification loss (Softmax Cross-Entropy) + bounding box regression loss (Smooth L1).
        \item Trained end-to-end in a two-stage pipeline:
        \begin{enumerate}
            \item Train RPN to generate high-quality proposals.
            \item Train classification and regression heads on the proposed regions.
        \end{enumerate}
    \end{itemize}
\end{itemize}

\subsection{DETR ( [DE]tection [TR]ansformer)}
\begin{itemize}
    \item \textbf{Backbone:}
        \begin{itemize}
            \item \textbf{ResNet-50:} Used as a feature extractor. It outputs a grid of image features that are passed into the transformer.
            \item \textbf{Positional Encoding:} Since transformers are permutation-invariant, positional encodings (sine and cosine based) are added to the backbone features to retain spatial information.
        \end{itemize}
        
    \item \textbf{Transformer Encoder-Decoder:}
        \begin{itemize}
            \item \textbf{Encoder:} A standard transformer encoder processes the flattened feature map. It consists of multi-head self-attention layers and feed-forward networks.
            \item \textbf{Decoder:} Uses a fixed number of learned object queries. Each query attends to the encoder output and attempts to predict an object.
            \item \textbf{Object Queries:} Learnable embeddings that allow the model to detect a fixed number of objects in a single forward pass.
        \end{itemize}
        
    \item \textbf{Prediction Head:}
        \begin{itemize}
            \item Consists of a shared feed-forward network applied to each decoder output.
            \item Each query outputs a bounding box and class label.
            \item DETR performs a bipartite matching (Hungarian algorithm) between predicted boxes and ground truth, enabling end-to-end optimization without anchor boxes or NMS.
        \end{itemize}

    \item \textbf{Inference and Evaluation:}
        \begin{itemize}
            \item Pretrained DETR (ResNet-50) from Hugging Face was used.
            \item Inference was performed on the Pascal VOC 2007 test set.
            \item Predictions were post-processed using the built-in DETR image processor with a confidence threshold of 0.7.
            \item Results were converted to COCO format and evaluated using \texttt{pycocotools} for standard object detection metrics (AP, AR).
        \end{itemize}
\end{itemize}

\vspace{-0.9}
\section{\textbf{Evaluation Metrics}}

\subsection{Intersection over Union (IoU)}
Intersection over Union (IoU) quantifies the degree of overlap between a predicted bounding box and the corresponding ground truth box. It is computed as:

\begin{equation}
\text{IoU} = \frac{\text{Area of Overlap}}{\text{Area of Union}}
\end{equation}

A higher IoU indicates better localization. Detection results are considered correct when their IoU with the ground truth exceeds a predefined threshold.

\subsection{Mean Average Precision (mAP)}
Mean Average Precision (mAP) is a standard metric used to evaluate the accuracy of object detection models. It combines both localization and classification performance. The mAP is computed as:

\begin{equation}
\text{mAP} = \frac{1}{N} \sum_{i=1}^{N} \text{AP}_i
\end{equation}

Here, \( N \) is the number of object categories, and \( \text{AP}_i \) is the Average Precision for the \( i \)-th class. In our experiments, we report:
\begin{itemize}
    \item \textbf{mAP@0.5:} Average Precision at a fixed IoU threshold of 0.5
    \item \textbf{mAP@0.5:0.95:} Average of AP values computed at multiple IoU thresholds from 0.5 to 0.95 in steps of 0.05
\end{itemize}


\section{Results}

\subsection{YOLO}

\subsubsection{COCO Dataset}
\begin{table}[H]
    \centering
    \caption{YOLO Performance on COCO Dataset}
    \begin{tabular}{lccc}
        \toprule
        \textbf{IoU Threshold} & \textbf{mAP} & \textbf{Precision} & \textbf{Recall} \\
        \midrule
        0.5:0.95 & 37.2 & 37.4 & 32.0 \\
        0.5      & 52.1 & 52.6 & 81.1 \\
        \bottomrule
    \end{tabular}
\end{table}

\subsubsection{Pascal VOC Dataset}
\begin{table}[H]
    \centering
    \caption{YOLO Performance on Pascal VOC Dataset}
    \begin{tabular}{lc}
        \toprule
        \textbf{Metric} & \textbf{Value} \\
        \midrule
        mAP@0.5        & 57.78  \\
        mAP@0.5:0.95   & 45.67  \\
        Precision      & 59.06  \\
        Recall         & 62.78  \\

        \bottomrule
    \end{tabular}
\end{table}

\subsection{Faster R-CNN}

\subsubsection{COCO Dataset}
\begin{table}[H]
    \centering
    \caption{Faster R-CNN Performance on COCO Dataset}
    \begin{tabular}{lccc}
        \toprule
        \textbf{IoU Threshold} & \textbf{mAP} & \textbf{Precision} & \textbf{Recall} \\
        \midrule
        0.5:0.95 & 34.44 & 37.8  & 42.23 \\
        0.5      & 53.83 & 46.19 & 52.13 \\
        \bottomrule
    \end{tabular}
\end{table}

\subsubsection{Pascal VOC Dataset}
\begin{table}[H]
    \centering
    \caption{Faster R-CNN Performance on Pascal VOC Dataset}
    \begin{tabular}{lc}
        \toprule
        \textbf{Metric} & \textbf{Value} \\
        \midrule
        mAP@0.5        & 73.35  \\
        mAP@0.5:0.95   & 46.387  \\
        avg\_IOU        & 0.511   \\
        Precision      & 17.77\\
        Recall         & 93.52 \\
        \bottomrule
    \end{tabular}
\end{table}

\subsection{DETR}

\subsubsection{COCO Dataset}
\begin{table}[H]
    \centering
    \caption{DETR Performance on COCO Dataset}
    \begin{tabular}{lccc}
        \toprule
        \textbf{IoU Threshold} & \textbf{mAP} & \textbf{Precision} & \textbf{Recall} \\
        \midrule
        0.5:0.95 & 36.5 & 39.0 & 42.1 \\
        0.5      & 52.3 & 46.5 & 49.7 \\
        \bottomrule
    \end{tabular}
\end{table}

\subsubsection{Pascal VOC Dataset}
\begin{table}[H]
    \centering
    \caption{DETR Performance on Pascal VOC Dataset}
    \begin{tabular}{lc}
        \toprule
        \textbf{Metric} & \textbf{Value} \\
        \midrule
        mAP@0.5        & 55.9  \\
        mAP@0.5:0.95   & 41.3  \\
        avg\_IOU       & 62.6    \\
        Precision      & 55.9    \\
        Recall         & 47.0  \\
        \bottomrule
    \end{tabular}
\end{table}

\section{Model Comparison and Analysis}

This section compares the performance of the three object detection models: YOLOv8, Faster R-CNN, and DETR—based on evaluation metrics computed on the COCO and Pascal VOC 2007 datasets.

\subsection{Performance Summary}

\subsubsection{COCO Dataset}

\begin{table}[H]
    \centering
    \caption{Performance Comparison on COCO Dataset}
    \begin{tabular}{lcccc}
        \toprule
        \textbf{Model} & \textbf{mAP@0.5} & \textbf{mAP@0.5:0.95} & \textbf{Precision} & \textbf{Recall} \\
        \midrule
        YOLOv8        & 52.1  & \textbf{37.2} & \textbf{52.6} & \textbf{81.1} \\
        Faster R-CNN  & \textbf{53.83} & 34.44 & 46.19 & 52.13 \\
        DETR          & 52.3  & 36.5  & 46.5  & 49.7 \\
        \bottomrule
    \end{tabular}
\end{table}

\subsubsection{Pascal VOC Dataset}

\begin{table}[H]
    \centering
    \caption{Performance Comparison on Pascal VOC 2007 Dataset}
    \begin{tabular}{lccccc}
        \toprule
        \textbf{Model} & \textbf{mAP@0.5} & \textbf{mAP@0.5:0.95} & \textbf{avg\_IoU} & \textbf{Precision} & \textbf{Recall} \\
        \midrule
        YOLOv8        & 57.78 & 45.67 & --     & \textbf{59.06} & 62.78 \\
        Faster R-CNN  & \textbf{73.35} & \textbf{46.39} & 51.1  & 17.77 & \textbf{93.52} \\
        DETR          & 55.9  & 41.3  & \textbf{62.6}  & 55.9  & 47.0 \\
        \bottomrule
    \end{tabular}
\end{table}

\subsection{Observations}

\begin{itemize}
    \item \textbf{YOLOv8} achieves the best recall and precision on the COCO dataset, making it ideal for real-time applications where fast and accurate detection is critical.
    
    \item \textbf{Faster R-CNN} outperforms the other models on the Pascal VOC dataset in terms of mean Average Precision (mAP@0.5), indicating strong performance in smaller, structured datasets. It also achieves the highest recall, although its precision is significantly lower, suggesting more false positives.
    
    \item \textbf{DETR} shows superior average IoU on Pascal VOC, indicating better bounding box alignment. It eliminates the need for anchor boxes and Non-Maximum Suppression (NMS), offering an end-to-end transformer-based detection pipeline.
\end{itemize}




\end{document}